\documentclass[aps,10pt,preprint,superscriptaddress,floatfix]{revtex4-2}

\usepackage{graphicx}
\usepackage{amsmath}
\usepackage{amssymb}
\usepackage{bm}
\usepackage[T1]{fontenc}
\usepackage{booktabs}
\usepackage{placeins}

\graphicspath{{figures/}}

\newcommand{\figorbox}[3]{%
  \IfFileExists{figures/#1.pdf}{\includegraphics[width=#3]{#1}}%
  {\IfFileExists{figures/#2.pdf}{\includegraphics[width=#3]{#2}}%
  {\fbox{\parbox[c][0.32\linewidth][c]{0.94#3}{\centering\ttfamily\small
   [figure \detokenize{#1} pending]}}}}%
}

\renewcommand{\thesection}{S\arabic{section}}
\renewcommand{\thetable}{S\arabic{table}}
\renewcommand{\thefigure}{S\arabic{figure}}
\renewcommand{\theequation}{S\arabic{equation}}

\begin{document}

\title{Supplemental Material for: Why sodium cobaltate superconducts only when it is wet}

\author{Santosh Kumar Radha}
% TODO confirm affiliation
\affiliation{TODO}

\author{Walter R. L. Lambrecht}
% TODO confirm
\affiliation{Department of Physics, Case Western Reserve University, Cleveland, Ohio 44106-7079, USA}

\date{\today}

\maketitle

This document is the engine room of the Letter. It carries the data, the derivations, and the
checks that the four-page format could only summarize, in the same order the Letter argues them:
methods, the well transition, the coupling model, the electronic evidence for the gallery band,
the water, the magnetism, the pairing state, the correlations we have not treated, the design rule
and its one counterexample, and the predictions that reach beyond this one compound.

\FloatBarrier
\section{Methods}
\label{sec:methods}

Every number quoted in the Letter comes from one of the calculation sets described here, in enough
detail that any of them can be checked or repeated without consulting the code. Table
\ref{tab:inventory} gives the inventory; the rest of this section gives the settings.

\emph{Code and functional.} All total energies are spin-polarized density-functional calculations
in the PBE generalized-gradient approximation, using projector-augmented-wave (PAW)
pseudopotentials from pslibrary, run in Quantum ESPRESSO 7.3.1 (GPU build). Plane-wave cutoffs are
60~Ry (wavefunctions) and 480~Ry (density) throughout. Electronic self-consistency uses
Marzari--Vanderbilt cold smearing of 0.02~Ry, local Thomas--Fermi charge mixing with mixing\_beta
$=0.3$, and a convergence threshold conv\_thr~$=10^{-7}$~Ry for all production total-energy scans;
the single exception is the Born--Oppenheimer molecular-dynamics stage (set~L,
Sec.~\ref{sec:water}), where conv\_thr is loosened to $10^{-6}$~Ry purely for throughput and every
other setting matches set~G. Every calculation is spin-polarized (nspin~$=2$); a direct
nspin~$=1$ vs.\ nspin~$=2$ check is reported below, and the local-spin-density over-polarization
that separately affects the $\sqrt3\times\sqrt3$ magnetic moments is discussed in
Sec.~\ref{sec:correlations}.

\emph{Cells.} The full-coverage $1\times1$ cell ($x=1$, one alkali per formula unit) is used for
the main displacement scans, in-plane lattice constants $a=2.888$~\AA{} (Na), $2.82$~\AA{} (Li),
and $2.90$~\AA{} (K; an ansatz upper estimate bracketing the experimental $2.84$--$2.87$~\AA{}
range for K$_x$CoO$_2$), with cell areas 7.22, 6.89, and $7.28$~\AA$^2$ respectively. The physical
$\sqrt3\times\sqrt3$ cell ($x=1/3$, $a=5.00$~\AA, area $21.67$~\AA$^2$) checks that the transition
survives at the real sodium content. Four CoO$_2$--CoO$_2$ spacings are computed for each series,
$c=5.5,\,6.9,\,8.4,\,9.9$~\AA, bracketing the anhydrous, monolayer-hydrate, and bilayer-hydrate
structures actually reported. The alkali sits on the Co-top in-plane site (the lowest-energy
registry found in the earlier surface work~\cite{RadhaLambrecht2021}) at $z=c/2+\delta$; the oxygen
height above the Co plane is held fixed at $z_O=0.96$~\AA{} for every spacing and every $\delta$,
the largest systematic in the calculation (below).

\emph{$k$-meshes.} $8\times8\times3$ for the $1\times1$ displacement scans, $6\times6\times3$ for
the $\sqrt3\times\sqrt3$ scans, $6\times6\times2$ (Gamma-centered) for the explicit-water cell
(set~G/J/L, Sec.~\ref{sec:water}), and a dense $24\times24\times8$ mesh for the
non-self-consistent eigenvalues that feed $N(0)$, the density of states, and the Fermi-surface maps
of Sec.~\ref{sec:fermi}. The band-structure/fatband runs of Sec.~\ref{sec:electronic} follow a
60-point $\Gamma$--M--K--$\Gamma$ path with {\sc projwfc.x} projections onto all Na $s/p$ and Co
$d$ orbitals.

\emph{Fitting and quantization.} $E(\delta)$ is fit to the even Landau polynomial
$E_0+\alpha\delta^2+\beta\delta^4$ (Eq.~\eqref{eq:landau}); where the quartic fit is not bracketed
by the sampled $\delta$ range (the $c\geq8.4$~\AA{} series, Sec.~\ref{sec:wells}), it is replaced
by a constrained sextic, $E_0+\alpha\delta^2+\beta\delta^4+\gamma_6\delta^6$ with $\gamma_6>0$
enforced and a hard-core cap $\delta_{\min}\leq c/2-2.0$~\AA{} (an ansatz for the alkali--oxygen
contact distance; varying it 1.8--2.5~\AA{} changes only the frozen-regime depths, no
conclusion). The quantized levels of Sec.~\ref{sec:softening} come from an exact one-dimensional
Schr\"odinger solve of the fitted well on a parity-resolved, staggered half-line finite-difference
grid ($N=2400$ points, box size following the well), which separates the even and odd sectors
exactly and resolves the near-degenerate tunnelling doublet without the numerical noise a
naive shooting method would add. The effective zero-point frequency is always reported as
$\hbar\omega_{\rm eff}=\hbar^2/2M\langle\delta^2\rangle_0$, never as $E_1-E_0$: in the deep
double-well regime $E_1-E_0$ is an exponentially small tunnelling splitting ($<10^{-12}$~meV for
Na at $c\geq6.9$~\AA), a real feature of the isolated double well but not the scale that couples to
electrons, whereas $E_2-E_0$, the even overtone quantified by $\langle\delta^2\rangle_0$, both sets
$\hbar\omega_{\rm eff}$ and is the boson of Eq.~\eqref{eq:lambda}.

\emph{Charge and coupling.} Bader charges are computed from the converged charge density (valence
counts $Z_{\rm val}=9$ for Na, 3 for Li, reflecting the semicore PAW pseudopotentials used). The
Allen--Dynes evaluation of Sec.~\ref{sec:model} uses $\mu^\ast=0.10$, is gated on $N(0)\geq0.1$
states~eV$^{-1}$cell$^{-1}$spin$^{-1}$ (results are insensitive to this threshold between 0.04 and
0.9), and is cut where $\lambda$ exceeds the polaron criterion $\lambda_{\rm loc}=2$ (surveyed range
1--2 in the literature), beyond which Allen--Dynes theory is not describing a well-defined
quasiparticle and raw numbers are not reported.

\emph{Water and its relaxation (sets G, J, L).} The rigid-cage scans (set~G) hold all twelve water
atoms fixed while the sodium scans $\delta$ (construction in Sec.~\ref{sec:water}). The adiabatic
scans (set~J) relax the same twelve water atoms by BFGS at each pinned sodium position, with
if\_pos freezing every Co, framework-oxygen, and sodium coordinate and leaving all water
coordinates free, force convergence threshold forc\_conv\_thr~$=10^{-3}$~Ry/Bohr and a cap of 100
ionic steps; the sixty-step-capped comparison run in Sec.~\ref{sec:water} uses the identical
settings, stopped early on purpose. The dispersion-corrected repeat (Sec.~\ref{sec:water}) adds
Grimme D3 (vdw\_corr~$=$~\texttt{'grimme-d3'}, supported natively in QE~7.3.1) on top of the set-J
settings. The Born--Oppenheimer molecular-dynamics stage (set~L) is generated with only the twelve
water atoms mobile (Verlet dynamics, stochastic-velocity-rescaling thermostat, $T=290$~K, time step
$\approx0.48$~fs, 4000 steps $\approx1.9$~ps); it is the open calculation referenced throughout this
work and has not yet been analyzed (Table~\ref{tab:inventory}, last row).

\emph{The frozen-oxygen systematic.} The largest systematic in the whole calculation is holding the
CoO$_2$ oxygen planes fixed at $z_O=0.96$~\AA. A direct check, moving the oxygen plane by
$\pm0.06$~\AA{} (to $z_O=0.90$ and $1.02$~\AA{} at $c=6.9$~\AA, $\delta=0.5$~\AA), shifts the
single-point energy by $+117$ and $-21$~meV respectively, comparable to the well depth itself; every
well depth quoted in this work therefore carries a systematic uncertainty of roughly $\pm50\%$. This
does not touch the sign of $\alpha$: the 5.5~\AA{} single well and the 6.9~\AA{} double well survive
the full range, and in the frozen polaronic regime the coupling is an order of magnitude past the
limit where a factor of two would matter. Two further checks bound the rest of the uncertainty
budget. A $12\times12\times4$ vs.\ $8\times8\times3$ $k$-mesh check at $c=6.9$~\AA, $\delta=0.5$~\AA{}
shifts the total energy by $-1.6$~meV, negligible on the $\sim100$~meV wells. An nspin~$=1$ vs.\
nspin~$=2$ check at the same geometry reproduces the well drop to within $0.2\%$ ($-89.1$ vs.\
$-88.9$~meV); at $c=9.9$~\AA, $\delta=0.75$~\AA{} the two spin treatments agree to $6.3\%$
($-121.5$ vs.\ $-113.8$~meV), so magnetism does not affect the well physics at a level that matters
for any conclusion in this work. Higher-cutoff (80~Ry) and lower-smearing (0.01~Ry) single points
exist at the reference geometry and are consistent with the production settings; the corresponding
relative (well-drop) checks are still running and are marked \% TODO verify where the manuscript
would otherwise imply they are complete. Composition dependence is checked directly rather than
assumed: interpolating $\alpha(c)$ for the physical $\sqrt3\times\sqrt3$ cell gives
$c^\ast\approx6.57$~\AA{} against $6.40$~\AA{} for the full-coverage $1\times1$ cell, a shift of
about $0.2$~\AA, so the transition is not an artifact of full coverage.

% ---- TABLE S1: calculation inventory ----
\begin{table*}[p]
  \caption{Calculation inventory. Every displacement scan, relaxation, and dense-mesh run behind
  the Letter and this Supplement, in the order it is used. ``Open'' marks the one calculation
  generated but not yet analyzed.}
  \label{tab:inventory}
  \scriptsize
  \renewcommand{\arraystretch}{1.35}%
  \begin{tabular}{@{}p{3.4cm}p{2.0cm}p{4.2cm}p{1.9cm}p{4.4cm}@{}}
    \toprule
    \textbf{Set / computes} & \textbf{Cell} & \textbf{Spacings / points} &
    \textbf{$k$-mesh} & \textbf{Feeds}\\
    \midrule
    A: anhydrous $E(\delta)$, Na/Li & $1\times1$ & $c=5.5,6.9,8.4,9.9$\,\AA;
      $\delta=0$--1.00\,\AA{} (6 pts) & $8\times8\times3$ &
      Fig.~S1, Fig.~2(a), Table~I rows 1,2,4\\
    B: anhydrous $E(\delta)$, Na & $\sqrt3\times\sqrt3$ & $c=5.5,6.9,8.4,9.9$\,\AA;
      $\delta=0$--0.75\,\AA{} (4 pts) & $6\times6\times3$ &
      Fig.~S1, the $\sqrt3$ check (Sec.~S1)\\
    C: dense NSCF + dos.x, $N(0)$ & both cells & same $c$ grid as A/B &
      $24\times24\times8$ & Fig.~2(b), Fig.~S3\\
    D: bands + {\sc projwfc} & $1\times1$, $\sqrt3\times\sqrt3$ &
      $c=5.5,6.9,9.9$\,\AA; $\Gamma$--M--K--$\Gamma$, 60 pts & --- &
      Fig.~3 (Letter), Fig.~S4\\
    E: K$_x$CoO$_2$ $E(\delta)$ & $1\times1$ & same $c$ grid as A &
      $8\times8\times3$ & Fig.~S9, $c^\ast_{\rm K}$\\
    F: gated LiCoO$_2$ (jellium) & $1\times1$ & $c=5.5$\,\AA; charge $0.1$--$0.3\,e$/f.u. &
      $8\times8\times3$ & gating prediction, Sec.~S12\\
    G: rigid-cage hydrate + vacuum ref. & $\sqrt3\times\sqrt3$ &
      $c=9.9$\,\AA; $\delta=0$--0.75\,\AA{} (4 pts) & $6\times6\times2$ &
      Fig.~4(a) (Letter)\\
    K: oxygen-height ($z_O$) scan & $\sqrt3\times\sqrt3$, vac. &
      $c=9.9$\,\AA; $z_O=0.90,1.02$\,\AA & $6\times6\times2$ &
      Fig.~S10, $\Delta v$ slope\\
    H: Na$_2$CoSe$_2$O check & primitive & SCF + dense NSCF &
      $24\times24\times8$ & negative control, Sec.~S12\\
    I2/J: adiabatic BFGS relax. & $\sqrt3\times\sqrt3$ &
      $c=9.9$\,\AA; $\delta=0$--1.00\,\AA{} (6 pts), 100 steps & $6\times6\times2$ &
      Fig.~4(a) (Letter)\\
    I3: DFT-D3 spot-check & $\sqrt3\times\sqrt3$ &
      $c=9.9$\,\AA; $\delta=0.00,0.15$ converged of 0.00--0.50 & $6\times6\times2$ &
      Sec.~S7 D3 check\\
    L/M: BOMD of set-G cell + snapshot rescans & $\sqrt3\times\sqrt3$ &
      $c=9.9$\,\AA; $\delta=0$, $T=290$\,K, $\sim$1.9\,ps & $6\times6\times2$ &
      \emph{open}: the decisive dynamical test\\
    \bottomrule
  \end{tabular}
\end{table*}

\FloatBarrier
\section{The well transition, in full}
\label{sec:wells}

This is the raw data behind the single number, $c^\ast_{\rm Na}\approx6.40$~\AA, that the Letter
builds on: the well shapes themselves, at every spacing and both cell sizes, and the check that the
transition is not an artifact of full sodium coverage.

For each spacing $c$ the sodium is displaced off centre by $\delta$ along $c$ and the total energy
$E(\delta)$ recorded; near the bottom the curve is even in $\delta$ and is fit by
\begin{equation}
  E(\delta) = E_0 + \alpha\,\delta^{2} + \beta\,\delta^{4},
  \label{eq:landau}
\end{equation}
continued to sixth order where the quartic term alone does not track the data. The sign of $\alpha$
is the whole story: while $\alpha>0$ the centre is a stable minimum and the sodium has one well;
once $\alpha<0$ the centre is a local maximum and the well has split in two. Figure~\ref{fig:wells}
shows the full set of curves for both cell sizes. At $c=5.5$~\AA{} (the dry spacing) $\alpha=+2.5$
eV/\AA$^2$ and the well is single and stiff; by $c=6.9$~\AA{} (the monolayer-hydrate spacing)
$\alpha=-0.5$~eV/\AA$^2$ and the ion sits in a pair of minima about 0.7~\AA{} off centre. The sign
flips between two spacings actually computed, so the crossing is bracketed,
$5.5~\text{\AA}<c^\ast_{\rm Na}<6.9$~\AA, and interpolating the four points places it at
$c^\ast_{\rm Na}\approx6.40$~\AA. Repeating the full scan at the physical composition $x=1/3$ in the
$\sqrt3\times\sqrt3$ cell moves the crossing by only about 0.2~\AA, so the transition is not an
artifact of the full-coverage cell.

Two of the curves in Fig.~\ref{fig:wells} (open symbols, $c\geq8.4$~\AA) are still falling at the
largest displacement sampled, meaning the sodium is sliding toward adsorption onto one CoO$_2$
sheet rather than settling into a genuine double well; the depths quoted for those spacings are
lower bounds, flagged as such in the figure, and this is exactly the pathology that the water cage
of Sec.~\ref{sec:water} is built to cure at $c=9.9$~\AA.

% ---- FIG S1: wells ----
\begin{figure}[t]
  \centering
  \figorbox{fig1_wells}{fig1}{0.62\columnwidth}
  \caption{The single-to-double-well transition, in full. Energy $E(\delta)$ of the off-centre
  sodium displacement, computed for Na $1\times1$ (top, $x=1$) and Na $\sqrt3\times\sqrt3$
  (bottom, $x=1/3$) at the four gallery spacings; points are total energies, curves the fits of
  Eq.~\eqref{eq:landau} continued to sixth order where needed. At 5.5~\AA{} the well is single and
  stiff ($\alpha>0$); by 6.9~\AA{} it has split into a double well ($\alpha<0$) with minima marked.
  Open symbols with arrows ($c\geq8.4$~\AA) flag scans still falling at the largest sampled
  displacement, where the sodium is adsorbing onto one layer and the quoted depth is a lower bound
  only.}
  \label{fig:wells}
\end{figure}

\FloatBarrier
\section{The two-band model and the coupling}
\label{sec:model}

The Letter states the pairing formula and defers the derivation here: why the linear coupling
vanishes, what the quadratic vertex is, and how $\lambda(c)$ and $\hbar\Omega(c)$ turn into the
dome of Fig.~2(c).

\emph{The bands.} The two states that share the gallery electron are borrowed, not fit here: an
SSH-type four-band tight-binding model developed for the surfaces of the layered
cobaltates~\cite{RadhaLambrecht2021}, restricted to the pair of bands relevant in the gallery, one
of alkali $sp_z$ character and one of CoO$_2$ character. Their band centres split according to how
the electron is shared between them, and that split is fixed by the Bader charge computed directly
in Sec.~\ref{sec:wells}, so the gallery band fills exactly as the two-dimensional gas turns on in
Fig.~2(b) of the Letter, with no free parameter left in the electronic part of the model. What is
added here is the sodium's own motion, and how the two bands respond to it. Figure~\ref{fig:model}
sketches all three pieces of the construction: the charge-sharing bands (panel a), the parity
argument (panel b), and the quantized double well with its overtone boson (panel c).

\emph{The parity argument.} The sodium displacement $\delta$ is odd under reflection through the
gallery midplane. A linear electron--phonon coupling would shift a band energy in proportion to
$\delta$, and such a term must vanish identically at the symmetric point by that reflection alone;
we checked this explicitly in the model at several $k_z$ and it holds everywhere in the gallery
band. The leading coupling is therefore quadratic, and the band energy responds as
$\tfrac12\chi\delta^2$, with
\begin{equation}
  \chi = 2\sum_m \frac{\lvert\langle {\rm alk}\lvert\,\partial H/\partial\delta\,\rvert m\rangle\rvert^2}
                       {E_{\rm alk}-E_m},
  \label{eq:chi}
\end{equation}
the ordinary second-order perturbation sum over the other bands $m$ of the model. Near the
transition this comes out to $\chi\approx-6$~eV/\AA$^2$, and it grows in magnitude as the band
splitting collapses further out in $c$.

\emph{The overtone boson.} Because the coupling is quadratic, an electron does not scatter off one
quantum of the rattling mode; it scatters off two. The matrix element that enters is
$\langle0\lvert\delta^2\rvert2\rangle$ between the ground state and the second excited state of the
quantized well (the first excited state is forbidden by the same parity that killed the linear
term), and the effective boson is the even overtone $\hbar\Omega=E_2-E_0$. The dimensionless
coupling is
\begin{equation}
  \lambda = \frac{2\,N(0)\,g^2}{\hbar\Omega}, \qquad
  g = \tfrac12\lvert\chi\rvert\,\langle0\lvert\delta^2\rvert2\rangle,
  \label{eq:lambda}
\end{equation}
with $N(0)$ taken directly from the density-functional calculation of Sec.~\ref{sec:wells}, not
adjusted to fit anything.

\emph{The dome and its two walls.} Feeding $\lambda(c)$, $\hbar\Omega(c)$, and $N(0)(c)$ into the
Allen--Dynes formula, with $\mu^\ast=0.10$ and gated on $N(0)\geq0.1$ states eV$^{-1}$cell$^{-1}$
spin$^{-1}$, gives the dome shown in Fig.~2(c) of the Letter: a narrow peak between $c\approx6.35$
and 6.45~\AA, height $T_c\approx14$~K at $\lambda\approx1.3$ and $\hbar\Omega\approx11$~meV. The two
walls are the two limits where the approximation itself tells us to stop trusting it. Below
$c^\ast$, $N(0)\to0$ and $\lambda\to0$ with it: there are no carriers to pair. Above the window, or
once the sodium is frozen into one well, $\hbar\Omega$ stays near 20~meV while $\lambda$ climbs past
2, which is no longer weak-to-intermediate coupling but a self-trapped polaron; Allen--Dynes theory
does not describe that regime, and a real bipolaron treatment past that cut has not been attempted
here. The ansatz range quoted for the peak, 11--23~K, comes from varying the electron--phonon range
parameter over its plausible values; the dome height depends on that parameter to the fourth power,
which is why the Letter states $T_c\approx14$~K as an order of magnitude and not a precision number.

% ---- FIG S2: the two-band coupling model ----
\begin{figure}[t]
  \centering
  \figorbox{fig3_model}{fig3_model}{\columnwidth}
  \caption{The two-band model and its coupling. (a) The gallery electron is shared between an
  alkali $sp_z$ band and a CoO$_2$ band, borrowed from the SSH-type surface
  model~\cite{RadhaLambrecht2021} and split by the Bader charge computed directly. (b) The gallery
  midplane mirror forbids a linear electron--phonon vertex at $\delta=0$; the leading coupling is
  quadratic. (c) The quantized double well, with the near-degenerate tunnelling doublet $E_0,E_1$
  and the even overtone $E_2$ that carries the quadratic coupling; the pairing boson is
  $\hbar\Omega=E_2-E_0$.}
  \label{fig:model}
\end{figure}

\FloatBarrier
\section{The Fermi surface and the ARPES comparison}
\label{sec:fermi}

The band structure of Sec.~\ref{sec:electronic} shows the gallery band crossing $E_F$; this section
checks what it leaves behind against the one photoemission measurement available for comparison.

Figure~\ref{fig:fermi} is the Fermi surface at $c=9.9$~\AA: the spectral function
$A(\mathbf{k},E_F)$ in the $k_z=0$ plane of the Na $1\times1$ hexagonal Brillouin zone, from the
dense $24\times24\times8$ non-self-consistent eigenvalues, symmetry-completed with the $C_{6v}$
point group and Lorentzian-broadened at $E_F$ (50~meV, both spins summed). The calculation gives a
single electron barrel around $\Gamma$, the sodium two-dimensional gas, and no $e_g'$ pockets at K.
Angle-resolved photoemission on the anhydrous parent Na$_{0.73}$CoO$_2$~\cite{Geck2007} sees the
same topology: a single $a_{1g}$ barrel around $\Gamma$ with measured $k_F\approx0.65$--0.70
\AA$^{-1}$, and no $e_g'$ pockets, the $e_g'$ band sitting $\approx85$~meV below $E_F$. Our barrel is
sodium- rather than cobalt-derived and the composition differs, so we compare topology and size, not
band identity; the two barrels are of comparable radius. The measured $a_{1g}$ velocities,
$v_F\approx0.3$ and 0.6~eV\,\AA{} along $\Gamma$--K and $\Gamma$--M respectively, run roughly a
factor of two below the band-theory value, a renormalization our self-energy-free Kohn--Sham map
does not carry (see also Sec.~\ref{sec:correlations}).

% ---- FIG S3: Fermi surface + ARPES comparison ----
\begin{figure}[t]
  \centering
  \figorbox{fig8}{fig8}{\columnwidth}
  \caption{Fermi surface at $c=9.9$~\AA{} and the ARPES comparison. $A(\mathbf{k},E_F)$ in the
  $k_z=0$ plane of the Na $1\times1$ hexagonal Brillouin zone, symmetry-completed and
  Lorentzian-broadened (50~meV, both spins summed): a single electron barrel around $\Gamma$ and no
  $e_g'$ pockets at K. Dashed: the ARPES barrel on anhydrous Na$_{0.73}$CoO$_2$~\cite{Geck2007}
  ($k_F\approx0.65$--0.70~\AA$^{-1}$), of comparable radius; the measured $a_{1g}$ velocities are
  $\sim2\times$ below the unrenormalized band value shown here.}
  \label{fig:fermi}
\end{figure}

\FloatBarrier
\section{The gallery band, and the electronic evidence for it}
\label{sec:electronic}

\subsection{At full sodium coverage}

Figure~\ref{fig:gallery} (reproduced from the Letter's Fig.~3, with the full caption repeated here
for completeness) is the band structure the Letter's headline number, 91\% sodium weight at $E_F$
by $c=9.9$~\AA, comes from. It is bare Kohn--Sham spectral weight, both spins summed and each band
Lorentzian-broadened by $\Gamma=40$~meV, with no self-energy applied. The cell is weakly
spin-polarized above the transition ($0.8\,\mu_B$ at $9.9$~\AA), so the gallery band is genuinely
spin-split; summing both spins is what a photoemission measurement would also collect.

\subsection{At the physical composition, $x=1/3$}
\label{sec:galleryfold}

The full-coverage result is checked at the real sodium content in the $\sqrt3\times\sqrt3$ cell.
Figure~\ref{fig:s3fold} mirrors the $x=1$ case exactly: no sodium weight at $E_F$ at 5.5~\AA, a
strong sodium band through $E_F$ at 9.9~\AA. These are supercell bands, and every one is a fold of
the primitive zone; lacking the plane-wave coefficients a weight-resolved unfolding would need, we
mark the folding relations rather than draw a misleading unfolded dispersion. The $1\times1$ K point
folds onto the supercell $\Gamma$, and the $1\times1$ M point folds onto the supercell M. Every band
in Fig.~\ref{fig:s3fold} should be read as folded.

% ---- FIG S4: sqrt3 folded bands ----
\begin{figure}[t]
  \centering
  \figorbox{fig9}{fig9}{\columnwidth}
  \caption{The gallery band at the physical composition, $x=1/3$. Kohn--Sham spectral weight of the
  Na $\sqrt3\times\sqrt3$ cell at (a) 5.5 and (b) 9.9~\AA, built exactly as in Fig.~\ref{fig:gallery}
  (grayscale total weight, blue sodium fatband, both spins summed, 40~meV broadening). No sodium
  weight sits at $E_F$ at 5.5~\AA; a strong sodium band crosses $E_F$ at 9.9~\AA, mirroring the
  $x=1$ result. The $1\times1$ K point folds onto the supercell $\Gamma$ and the $1\times1$ M point
  onto the supercell M; every band shown here is folded.}
  \label{fig:s3fold}
\end{figure}

\begin{figure*}[t]
  \centering
  \figorbox{fig7}{fig7}{0.98\textwidth}
  \caption{(Reproduced from the Letter.) The gallery band appears. Kohn--Sham spectral weight of Na
  $1\times1$ along $\Gamma$--M--K--$\Gamma$ at (a) 5.5, (b) 6.9, and (c) 9.9~\AA, referenced to
  $E_F$ (dashed). Grayscale is the total DFT spectral weight (all bands, both spins summed, each
  Lorentzian-broadened by $\Gamma=40$~meV); the blue fatband is the sodium-orbital projected weight
  from {\sc projwfc}, its width proportional to sodium character (key in panel c). At 5.5~\AA{} the
  $x=1$ cell is a band insulator and the sodium band lies $\approx2$~eV above $E_F$; by 9.9~\AA{} a
  band of $\sim$91\% sodium character has descended through $E_F$, the interlayer electron gas made
  visible as a single band. No self-energy is applied: this is bare Kohn--Sham weight, not a
  renormalized quasiparticle spectrum.}
  \label{fig:gallery}
\end{figure*}

\FloatBarrier
\section{Mode softening and the quantized wells}
\label{sec:softening}

The pairing model of Sec.~\ref{sec:model} needs a soft anharmonic mode, not just a sign change in
$\alpha$; this section shows how soft it actually gets.

Figure~\ref{fig:soft} is the effective zero-point frequency $\hbar\omega_{\rm eff}=\hbar^2/2M
\langle\delta^2\rangle_0$ of the quantized alkali mode, obtained from an exact one-dimensional
solution of the fitted wells of Sec.~\ref{sec:wells} at the physical sodium mass. It collapses by
three orders of magnitude across the transition, from a stiff $\sim30$~meV single-well mode to a
sub-meV soft mode as the gallery opens, entering a quantum-paraelectric regime in which the ion is
smeared across the barrier rather than localized in either minimum. The inset shows the lowest
levels of the 6.9~\AA{} sodium double well directly: the near-degenerate tunnelling doublet
$E_0,E_1$ that the parity argument of Sec.~\ref{sec:model} forbids from coupling linearly, and the
even overtone $E_2$ that carries the quadratic vertex and sets $\hbar\Omega$.

% ---- FIG S5: softening ----
\begin{figure}[t]
  \centering
  \figorbox{fig5_softening}{fig5}{\columnwidth}
  \caption{The sodium mode softening. Effective zero-point frequency
  $\hbar\omega_{\rm eff}=\hbar^2/2M\langle\delta^2\rangle_0$ of the quantized alkali mode, from an
  exact one-dimensional solution of the fitted wells at the physical mass. It collapses by three
  orders of magnitude across the transition, from a stiff $\sim30$~meV single-well mode to a
  sub-meV soft mode, entering a quantum-paraelectric window. Inset: the lowest levels of the
  6.9~\AA{} double well, the tunnelling doublet $E_0,E_1$ and the overtone $E_2$ that carries the
  quadratic coupling.}
  \label{fig:soft}
\end{figure}

\FloatBarrier
\section{The water cage}
\label{sec:water}

The Letter's Fig.~4 rests entirely on the construction described here. We give it in full because
the result turns on the construction, and because the three-way comparison (vacuum, rigid cage,
adiabatic cage) is what turns a plausible story into a checkable calculation.

\subsection{Construction}

The host is the $\sqrt3\times\sqrt3$ ($x=1/3$) Na$_x$CoO$_2$ cell at the bilayer spacing
$c=9.9$~\AA{} ($a=5.00$~\AA), the same cell and the same spin-polarized PBE+PAW settings of
Sec.~\ref{sec:methods} used for the anhydrous scans. The CoO$_2$ sheets and oxygen heights are
frozen at the values used throughout, so the water is the only thing added. Each water is a rigid
gas-phase molecule, O--H bond 0.9572~\AA{} and H--O--H angle $104.52^\circ$. Four are placed per
cell, in the arrangement neutron diffraction finds for the bilayer
hydrate~\cite{Jorgensen2003,Lynn2003}: the sodium sits on the Co-top $(0,0)$ site on its plane at
$z=c/2+\delta$, and the four water oxygens lie in two layers 1.2~\AA{} below and above that plane,
on in-plane sites the sodium does not occupy, two O-top images in the lower layer and two Co-top
images in the upper layer. Over the whole displacement scan the sodium--water oxygen distances run
2.05--2.57~\AA{} (lower layer) and 2.92--3.13~\AA{} (upper), an octahedral-like first coordination
shell. The hydrogens point away from the sodium, toward the nearest CoO$_2$ sheet and radially
outward, the molecular bisector tilted $50^\circ$ off the $c$ axis, chosen to maximize the smallest
intermolecular contact: every water--water, water--framework, and water--sodium distance stays
$\geq2.05$~\AA{} across the whole $\delta$ grid. The two down-pointing and two up-pointing molecules
cancel the layer's $z$-dipole exactly, so the cage adds no net dipole field.

\subsection{The rigid cage (set G)}

The waters are held fixed while the sodium scans $\delta$, so $E_{\rm hyd}(\delta)-E_{\rm
vac}(\delta)$, the hydrate curve minus the identical cell with the twelve water atoms deleted, is
the clean effect of the solvation cage on the sodium potential, with no relaxation confounds. This
probes the static cage only: it answers whether an expanded gallery leaves the sodium chemisorbed
(it does, in vacuum, an unstable mode with $\lvert\hbar\omega\rvert\approx11i$~meV) or bound in a
confined well (it does, with the cage, $\hbar\omega_{\rm eff}\approx23$~meV), and fixes that well's
frequency. It cannot represent the water's own motion.

\subsection{Adiabatic relaxation (set J) and its surprise}

The obvious next move is to let the cage relax. We repeated the scan with the four waters no longer
frozen but allowed to find their own minimum-energy configuration at each pinned sodium position, a
full BFGS relaxation of the twelve water atoms, one hundred ionic steps, converged at every
$\delta$. A relaxing shell, if water were a lubricant in the naive sense, would slide out of the
sodium's way and flatten the well; instead it follows the sodium off-centre and deepens the bound
well, from 199~meV (frozen) to 380~meV at the same $\delta\approx0.30$~\AA{} minimum, and quantized
about that minimum the adiabatic mode is stiffer, not softer, near 37~meV. Classical adiabatic water
does not soften the sodium; taken at face value, the naive lubricant picture is wrong. A
dispersion-corrected (DFT-D3) repeat tracks this deepening at the two points where the D3
trajectory has converged, $-210$~meV against $-222$~meV (PBE) at $\delta=0.15$~\AA; the remaining D3
points are still relaxing but are consistent with the same trend, so the deepening is not a
van-der-Waals artifact of the bare PBE cage.

\subsection{The 0.4 eV spread, and the sixty-step artifact}

The single relaxed configuration is only the deepest point of a landscape the four-water
hydrogen-bond network does not sit still in. A relaxation of the same $\delta=0.30$~\AA{} cell,
stopped at sixty ionic steps rather than carried to full convergence, comes to rest 394~meV above
the fully converged minimum, a different, shallower local minimum of the cage at the identical
sodium position (the open star in the Letter's Fig.~4(a)). The hydrogen-bond network therefore
presents the sodium not with one potential but with a family of them, spanning at least 0.4~eV
between local minima at fixed $\delta$; the sodium potential is set by the instantaneous cage
configuration, of which the deep adiabatic surface is only the lower envelope. (This is also, for
the record, why an earlier reading of the sixty-step run as a near-flat 5~meV well was an artifact
of incomplete relaxation, not a physical softening; the corrected, fully converged comparison is the
one reported throughout.)

\subsection{The timescale argument}

The two motions live on different clocks, and that is what closes the argument. The bound sodium
mode has $\hbar\omega\approx23$~meV, a period near 200~fs; the hydrogen-bond network of gallery
water reorganizes on the picosecond scale that quasielastic neutron scattering measures
directly~\cite{Jalarvo2008}. The sodium is fast against the cage. It does not ride the deep
adiabatic surface, which the slow network can only present once fully relaxed; it rattles on a
rigid-cage-like surface drawn, at each instant, from the thermal and quantum ensemble of cage
configurations, an ensemble that, as the 0.4~eV spread shows, includes surfaces far shallower than
the adiabatic minimum. The softening this mechanism needs is therefore statistical-dynamical, not
adiabatic: the fast sodium samples the soft members of the cage ensemble, it does not wait for the
slow cage to relax to its own stiff, deep minimum. The frozen-cage well is one member of that
ensemble and the adiabatic well its extreme deep member; the physical mode lives between them, and
the one calculation that can show this directly, a finite-temperature Born--Oppenheimer or
path-integral molecular-dynamics sampling of the coupled sodium--water system, remains open.

\subsection{The isotope-null consistency check}

This picture makes the deuteration and $^{18}$O experiments consistency checks rather than
predictions, because both have already been run. Replacing H$_2$O by D$_2$O shifts $T_c$ by less
than 0.2~K~\cite{Jin2003}; replacing $^{16}$O by $^{18}$O is also null~\cite{Yokoi2008}, with the
authors of that study noting explicitly that a null shift does not by itself exclude a phonon
mechanism once strong correlations are present. Neither result is evidence against the picture
argued here, because neither isotope substitution touches the boson: the pairing mode is the sodium
rattling along $c$, gated by the statistical ensemble of cage configurations, and the water's own
mass enters only through how stiff the hydrogen-bond network is, a second-order effect on the
ensemble the sodium samples rather than a first-order mass term in the boson frequency itself. A
mechanism that paired electrons directly off a water vibration, by contrast, predicts a
first-order isotope shift and does not survive this result. Sodium itself offers no equivalent test:
$^{23}$Na is its only stable isotope, so the direct mass test on the boson is unavailable, which is
why Sec.~\ref{sec:predictions} emphasizes the spectroscopic search for the mode itself over any
further isotope substitution.

\FloatBarrier
\section{Spin phase diagram and spin-orbit coupling}
\label{sec:spin}

The gas that condenses in the gallery is not always paramagnetic, and where it orders matters for
two old puzzles in this compound: why a magnetic phase sits next to the dome, and why the
Knight-shift literature reports both a singlet and a triplet. This section builds the phase diagram
that answers both.

\subsection{The Stoner strip}

Right where the gas first condenses, the gallery band is narrow and its density of states high, and
a narrow, high-density-of-states band is a Stoner magnet. The Radha--Lambrecht surface calculation
already finds this polarized gallery gas~\cite{RadhaLambrecht2021}, with the gallery moment
antiparallel to the cobalt. Carrying the same SSH-type surface-charge model into the gallery and
applying a Stoner criterion to the gallery band gives, in the plane of sodium content $x$ and
CoO$_2$ spacing $c$, a polarized strip that hugs the carrier turn-on line, boosted to a
paramagnon-rich metal just outside it and settling into a plain paramagnet beyond
(Fig.~\ref{fig:spinpd}(a)); the spacing dependence of the enhancement is the cut in
Fig.~\ref{fig:spinpd}(b). The diagram is anchored to our own numbers: the density-functional moment
switches on exactly at the well transition (Fig.~\ref{fig:spinpd}(c)), which is where the model
places the carrier turn-on at full coverage. The one number we do not trust is the moment of the
$\sqrt3\times\sqrt3$ cell, over-polarized in the local-spin-density approximation and excluded from
the diagram; the well energies on which the mechanism rests shift by less than 6\% under the spin
treatment (Sec.~\ref{sec:correlations}), so the exclusion leaves the structural story intact.

In this picture the magnetism is a thin strip hugging the carrier turn-on, strongest just past it
and dying as the band widens, so superconductivity lives one step beyond the magnetism and the
magnetism lives one step beyond the turn-on. This is the ordering Sakurai, Ihara, and Takada report
for the magnetic phase bordering the dome against cobalt valence~\cite{Sakurai2015}. Figure
\ref{fig:dome} puts real data on that claim: the experimental $T_c$ against cobalt valence, with the
superconducting dome and the adjacent magnetic-phase region shown as pastel fills, is the anchor
for ``magnetism borders the dome'' rather than a schematic restatement of it. Proximity to the
Stoner boundary also dissolves an old contradiction. Early powder NMR suggested a triplet, later
single-crystal work a singlet; in this picture that is not a contradiction but a statement of where
each sample sits, because the Stoner enhancement varies steeply across the strip and hydration
drifts from sample to sample, and right at the boundary an equal-spin channel of $f$ symmetry lies
close in energy to the singlet, the link to the two surviving $f$-wave states of
Sec.~\ref{sec:gauntlet} that Mazin and Johannes could not eliminate~\cite{MazinJohannes2005}. The
prediction is sharp: the NMR Knight-shift suppression below $T_c$ should anticorrelate, across
samples, with the normal-state Stoner enhancement. Both symmetry-allowed states of Mazin and
Johannes are consistent with the muon null~\cite{Higemoto2004}, and we do not need the chiral
$d+id$ scenario~\cite{Kiesel2013} to fit the record, though our data do not exclude it.

% ---- FIG S6: gallery spin phase diagram (three panels) ----
\begin{figure*}[t]
  \centering
  \figorbox{fig10}{fig10}{0.92\textwidth}
  \caption{The gallery-2DEG spin phase diagram. (a) Stoner phase diagram of the interlayer gas,
  built by carrying the SSH-type surface-charge model of Ref.~\cite{RadhaLambrecht2021} into the
  gallery: in the plane of sodium content $x$ and CoO$_2$ spacing $c$, the gallery band is empty
  below the carrier turn-on line, narrow and Stoner-polarized (red) just above it, then a
  paramagnon-boosted (amber) and plain paramagnetic (blue) metal further out. Markers: the
  monolayer (6.9~\AA, $\times$) and bilayer (9.9~\AA, $\circ$) hydrates at the physical $x=0.35$.
  (b) Cuts at $x=0.35$: the carrier fraction $q$ per Na and the Stoner enhancement $S$ against $c$.
  (c) The DFT cell-integrated moment $\langle\lvert m\rvert\rangle$ for Na and Li $1\times1$,
  switching on at the same spacing the well transition does; the $\sqrt3\times\sqrt3$ cell is
  over-polarized in LSDA and shown only as an excluded caveat band. The model is an ansatz
  calibrated to the one on/off hydrate anchor pair; every parameter is source-tagged in the
  generating script.}
  \label{fig:spinpd}
\end{figure*}

% ---- FIG S7: Tc vs Co valence with SC dome + magnetic region ----
\begin{figure}[t]
  \centering
  \figorbox{figS_dome}{figS_dome}{\columnwidth}
  \caption{Magnetism borders the dome, against real data. Experimental $T_c$ against the
  redox-\emph{titrated} cobalt valence, the eight batches of Milne \textit{et al.}~\cite{Milne2004},
  with the superconducting dome and the adjacent magnetic-phase region drawn as pastel fills; the
  nominal-composition ($4-x$) points of Ref.~\cite{Schaak2003} use a different, offset ($\approx
  0.2$--$0.3$) valence axis and are deliberately omitted rather than mixed with the titrated set.
  The magnetic-phase boundary ($v\approx3.36$) is anchored to the collapse threshold Milne \textit{et
  al.} report, combined with the phase adjacency Sakurai \textit{et al.} give~\cite{Sakurai2015}, and
  is approximate. The titrated batches cover the optimal plateau and the overdoped collapse, not the
  underdoped rising edge, so this is a partial map of the dome, not a claim of its full shape; what it
  shows is real: the magnetic region sits immediately outside the superconducting one, the adjacency
  the Stoner strip of Fig.~\ref{fig:spinpd}(a) predicts from the calculation alone.}
  \label{fig:dome}
\end{figure}

\subsection{Spin-orbit coupling}

Spin-orbit coupling does not select or mix the pairing at any level that matters here. The gallery
state is sodium-derived, with only a few percent cobalt weight, so its effective on-band spin-orbit
scale is a couple of meV; the Rashba splitting from the off-centre sodium, estimated from the
displaced charge, is about 0.8~meV, of order the gap but far below the Fermi energy, and the
parity-mixed component of the gap it induces is at the sub-percent level. On the average
centrosymmetric structure the up and down sectors carry opposite Rashba fields and largely cancel.
Spin-orbit coupling is a small correction to the spin response, not a term in the pairing.

\FloatBarrier
\section{The pairing gauntlet}
\label{sec:gauntlet}

Twenty years of proposals for Na$_x$CoO$_2\cdot y$H$_2$O have been thinned by four measurements;
Table~\ref{tab:gauntlet} sets the surviving field, including the state of this work, against them,
in a form where no verdict runs past two printed lines.

The natural ground state of a phonon-mediated instability on the gallery gas is a spin singlet,
anisotropic and $s_\pm$-like across the two bands, the gallery two-dimensional gas and the $a_{1g}$
pocket, and it breaks no time-reversal symmetry. That single fact clears the two constraints that do
the most damage in Table~\ref{tab:gauntlet}: the singlet passes the single-crystal Knight shift that
kills the triplet proposals, and the preserved time-reversal symmetry passes the $\mu$SR
null~\cite{Higemoto2004} that kills the chiral $d+id$ states~\cite{Kiesel2013}.

The other two constraints our picture turns to account rather than merely surviving. There is no
Hebel--Slichter peak because the pairing boson is a soft, strongly anharmonic mode, and the
anharmonic broadening that softens it also smears the coherence peak; we state this as expected of a
soft-mode coupling, not as a fitted success. The specific-heat nodes are the more telling case. The
gap on the gallery band is small, set by the dilute gallery density of states, and a small gap on a
two-dimensional gas is fragile against precisely the gallery disorder that dehydration introduces. A
fragile second gap driven below the measurement threshold looks like line nodes, which explains both
the node claims and their notorious dependence on sample age~\cite{Sakurai2015}: the gap is not
nodal, it carries a soft second component the aging of the water network quietly removes. The state
is falsifiable in two clean ways: any breaking of time-reversal symmetry, a $\mu$SR relaxation onset
at $T_c$ or a polar Kerr rotation, excludes the singlet channel, and so does a Knight shift that
stays flat below $T_c$ in a fresh single crystal, because a singlet must show the shift drop. Either
observation would end the picture.

% ---- TABLE S1: pairing gauntlet ----
\begin{table*}[t]
  \caption{The pairing gauntlet. Verdicts are single words read against each state's own symmetry:
  pass, fail, or tension (a constraint met only with an added assumption, flagged
  $^{\rm a\text{--}f}$, notes below). The Sakurai entry is a mechanism, not a fixed symmetry, so the
  symmetry columns do not constrain it. The last row is the state argued in the Letter
  (Sec.~\ref{sec:gauntlet}).}
  \label{tab:gauntlet}
  \renewcommand{\arraystretch}{1.3}%
  \begin{tabular}{@{}p{4.4cm}p{2.0cm}p{2.0cm}p{2.0cm}p{2.0cm}@{}}
    \toprule
    \textbf{Proposed state} & \textbf{Knight shift} & \textbf{Hebel--Slichter} &
    \textbf{Heat nodes} & \textbf{$\mu$SR TRS}\\
    \midrule
    RVB singlet $d+id$~\cite{Baskaran2003,KumarShastry2003,WangLeeLee2004}
      & pass & pass & tension$^{\rm a}$ & fail\\
    Spin-triplet $p$/$f$~\cite{TanakaHu2003,KurokiTanakaArita2004}
      & fail & pass & pass & pass$^{\rm b}$\\
    Chiral $d+id$, fRG~\cite{Kiesel2013}
      & pass & tension$^{\rm a}$ & tension$^{\rm a}$ & fail\\
    Two surviving $f$-wave~\cite{MazinJohannes2005}
      & tension$^{\rm c}$ & pass & pass & pass\\
    $s$-wave, anharmonic H/O phonons~\cite{Lynn2003}
      & pass & fail$^{\rm d}$ & fail$^{\rm e}$ & pass\\
    Magnetic-fluctuation~\cite{Sakurai2015} & --- & --- & --- & ---\\
    \midrule
    \textbf{This work}: singlet $s_\pm$, TRS-preserving, two-gap
      & pass & pass$^{\rm f}$ & pass$^{\rm g}$ & pass\\
    \bottomrule
  \end{tabular}\\[4pt]
  \parbox{0.98\textwidth}{\footnotesize
  $^{\rm a}$~Fully gapped states must invoke disorder or a second gap to explain the observed
  $T^2$ term.
  $^{\rm b}$~Requires $e_g'$ pockets that ARPES does not see~\cite{Geck2007}.
  $^{\rm c}$~Passes only after an added reconciliation of the singlet-favoring Knight shift with a
  triplet order parameter.
  $^{\rm d}$~Predicts a Hebel--Slichter peak the data do not show.
  $^{\rm e}$~Fully gapped; does not reproduce the observed low-$T$ nodal-like term.
  $^{\rm f}$~Anharmonic broadening of the soft-mode boson smears the coherence peak.
  $^{\rm g}$~Small gallery gap is disorder-fragile; explains both the node claims and their
  sample-age dependence.}
\end{table*}

\FloatBarrier
\section{Correlations we have not treated}
\label{sec:correlations}

Every estimate above is built on Kohn--Sham bands, and the cobaltates are correlated; this section
states which way that correlation pushes the numbers, rather than leaving the reader to guess.

Angle-resolved photoemission on the anhydrous parent finds dispersions roughly twice flatter than
LDA, with $a_{1g}$ velocities $v_F\approx0.3$--0.6~eV\,\AA{} against a band value about twice
larger~\cite{Geck2007}. A mass renormalization $m^\ast/m\approx2$ raises the density of states
$N(0)$ and lowers $v_F$ in the same proportion, and $N(0)$ enters the coupling in the direction that
increases it: since $\lambda=2N(0)g^2/\hbar\Omega$ was built on the unrenormalized
density-functional $N(0)$, the $\lambda$ and $T_c$ estimates of Sec.~\ref{sec:model} are, if
anything, conservative. The same flattening deepens the narrow-band limit the gallery gas already
sits in, and a flatter band at fixed filling carries a higher $N(0)$, which only sharpens the
Stoner-strip argument of Sec.~\ref{sec:spin}, correlations widen the polarized region of
Fig.~\ref{fig:spinpd} rather than narrow it. The natural next step is a DFT+DMFT calculation on the
gallery band, which we have not attempted. One caution runs the other way: correlations do not only
renormalize the band, they can also broaden and damp the sodium mode that carries the pairing, and a
strongly overdamped mode would couple less efficiently than the sharp Einstein boson assumed here.

The rest of the list is shorter. The transition spacing is computed at full sodium coverage, $x=1$,
and checked at the physical $x=1/3$ in a single cell, where it moves by about 0.2~\AA. The
local-spin-density treatment over-polarizes the $\sqrt3$ cell, giving an unphysically large and flat
moment there; we rest the magnetic story on the $1\times1$ series and exclude the $\sqrt3$ moments,
and the well energies themselves are insensitive to the spin treatment at the few-percent level, so
the over-polarization does not touch the wells. The explicit water (Sec.~\ref{sec:water}) is treated
statically, a rigid cage frozen while the sodium scans and an adiabatically relaxed cage at each
sodium position; together they settle the static question and show that adiabatic relaxation
deepens rather than softens that well, but neither can represent the statistical-dynamical softening
the moving network supplies, the one calculation left open. The pairing is treated as a single
Einstein boson in the Allen--Dynes approximation, with an explicit cut where the coupling leaves the
regime that approximation can describe; a real bipolaron treatment past that cut is not attempted.

\FloatBarrier
\section{The design rule, prior art, and the LiZrNCl null hypothesis}
\label{sec:designrule}

If the picture holds, it is not just an account of one compound, it is a design rule; this section
states the rule, credits the ingredients that are not new, and then tests it against the one
layered intercalate that seems, at first look, to break it.

Every ingredient of the rule is old. Pairing tuned to a soft polar mode at a quantum-critical point
is the SrTiO$_3$ story~\cite{Edge2015,Muller1979QuantumParaelectric,vanderMarel2019SrTiO3,
Yu2022QuantumParaelectric}; an occupied interlayer state as the switch for superconductivity in an
intercalate, set by the layer separation, is the graphite-intercalation
story~\cite{Csanyi2005,OhtomoHwang2004}; enhancement of $T_c$ as an intercalant mode softens, then
collapse when the ion orders, is seen in CaC$_6$ under pressure~\cite{Gauzzi2007}; a rattling alkali
ion in an oversized cage driving strong-coupling pairing is the $\beta$-pyrochlore
story~\cite{Yonezawa2004KOs2O6,Hiroi2012RattlingReview}; and carriers in one layer paired by a
quantum-paraelectric layer next door is an old model of Shenoy, Subrahmanyam, and
Bishop~\cite{Shenoy1997}. What is new here is the joinery: the donor ion is itself the
quantum-paraelectric unit, not a caged spectator or a host-lattice mode, and one structural knob,
the gallery spacing, turns on the carriers and tunes the ion through its bistability at once, so the
window closes on both sides rather than one. Stated as a rule: open a layered gallery past the
spacing where the donor ion turns bistable, and the same motion that spills the ion's electrons
into an interlayer gas softens the ion into the pairing boson.

A rule stated that broadly invites a counterexample, and the most-studied layered donor intercalate
is one. In Li$_x$ZrNCl and its hafnium analogue~\cite{Yamanaka1998}, opening the gallery from 9 to
more than 22~\AA{} never kills the superconductivity: $T_c$ rises modestly and then saturates
(Fig.~\ref{fig:zrncl}(a)), and it survives with no intercalant ion at all, whether the chlorine is
de-intercalated or the carriers are supplied by a gate~\cite{Nakagawa2021}. This looks like a
refutation until the premise is checked, and the premise fails: here the donor's electron does not
return to the gallery. It fills an in-plane Zr/Hf $d$--N $p$ host band, the ion is chemically locked
and spectroscopically silent, and the pairing belongs to the host, not to any gallery gas or soft
ion. The nitride is therefore orthogonal to our rule, it tests a different pairing channel, and its
use to us is disciplinary. It forbids the universal statement: the rule holds only where the
interlayer state lies at the Fermi level, the precondition made computable by Cs\'anyi and
co-workers~\cite{Csanyi2005}, and fails where the host band lies below it. And it fixes the null
hypothesis. A gallery opened without a soft donor ion gives a monotonic, saturating $T_c(c)$; the
nonmonotonic dome computed for the hydrate (Fig.~\ref{fig:zrncl}(b)) is, against that background,
the anomaly that demands the extra ingredient, the soft sodium, that the nitride does not have.

% ---- FIG S8: ZrNCl comparison ----
\begin{figure}[t]
  \centering
  \figorbox{figS_zrncl}{figS_zrncl}{\columnwidth}
  \caption{The null hypothesis, tested. (a) Literature $T_c$ against gallery spacing for
  Li$_x$ZrNCl and its hafnium analogue~\cite{Yamanaka1998,Nakagawa2021} (17 points spanning both
  chemical and gate-induced intercalation): opening the gallery from 9~\AA{} to beyond 22~\AA{}
  gives a monotonic, saturating $T_c(c)$, because the donor electron never returns to the gallery.
  (b) Our computed Allen--Dynes dome for the hydrate, with the 4.5~K bilayer-hydrate anchor (star),
  rises and falls instead: a nonmonotonic window that appears only when the gallery hosts a soft
  donor ion. The comparison in (a) and (b) is the null hypothesis this work must beat.}
  \label{fig:zrncl}
\end{figure}

\FloatBarrier
\section{Predictions beyond the compound}
\label{sec:predictions}

If the rule of Sec.~\ref{sec:designrule} is right, it should say something about the neighboring
compounds and about probes that have not yet been tried on this one; this section collects those
predictions and, where we could, puts a number on them.

\subsection{Lithium, sodium, and potassium}

Lithium cobaltate is the same structure with a smaller ion, and it has never been made to
superconduct; the picture requires that. Its well crosses into the double-well regime only at the
smallest spacing in the dataset, so there is only an upper bound $c^\ast_{\rm Li}\leq5.5$~\AA, and
at that spacing the well is a mere 4~meV deep while the lithium zero-point energy is 8.8~meV,
larger than the well itself. The light ion does not sit in either minimum; it is smeared across
both, a quantum paraelectric, dynamically centred despite the negative $\alpha$, and a centred ion
is a stiff ion with no soft anharmonic mode to offer. Lithium fails on the glue, for the plain
reason that it is light. A heavier gallery ion should push the threshold the other way: for
potassium the Landau coefficient crosses zero at $c^\ast_{\rm K}\approx6.75$~\AA, against 6.40~\AA{}
for sodium and $\leq5.5$~\AA{} for lithium, ordering Li~$<$~Na~$<$~K with ionic size
(Fig.~\ref{fig:thresholds}); hydrated to that spacing, potassium cobaltate should superconduct by
the same route with a heavier, softer mode, a clean prediction for hydrated K$_x$CoO$_2$.

% ---- FIG S9: Li/Na/K thresholds ----
\begin{figure}[t]
  \centering
  \figorbox{figS_thresholds}{figS_thresholds}{\columnwidth}
  \caption{The threshold ordering, in one plot. $\alpha(c)$ for Li, Na, and K $1\times1$, with the
  three thresholds marked: $c^\ast_{\rm Li}\leq5.5$~\AA{} (an upper bound only, $\alpha$ is already
  negative at the smallest spacing sampled, not a bracketed crossing), $c^\ast_{\rm
  Na}\approx6.40$~\AA, and $c^\ast_{\rm K}\approx6.75$~\AA{} (both bracketed crossings). The
  threshold rises with ionic size, the prediction behind the hydrated-K$_x$CoO$_2$ test proposed in
  the Letter.}
  \label{fig:thresholds}
\end{figure}

\subsection{The valence fingerprint}

The gallery gas leaves a fingerprint on the sodium valence that two kinds of probe should read
differently. If the sodium keeps a fraction $q$ of its electron in the gallery rather than giving it
to the cobalt, a site-specific valence probe, the cobalt $L$-edge or NQR, and a charge-counting
probe, redox titration of the whole crystal, should disagree, and only once the gallery is open. We
can put a number on it: the Bader charge on the sodium in the $x=1/3$ cell rises from the
empty-gallery monolayer spacing to the superconducting bilayer by $\Delta q\approx0.18\,e$, so the
gallery-specific valence split is $\Delta v=x\,\Delta q\approx0.06$ per cobalt, switching on with
the water and vanishing in the monolayer hydrate where the gallery band is empty. This small split
is not the large valence offset seen in these materials~\cite{Milne2004}, which needs an extra
electron reservoir such as oxonium and is several times larger; ours is the small, gallery-specific
piece alone. Figure~\ref{fig:zscan} gives the microscopic slope behind the number: the Landau
coefficient $\alpha$ against the frozen oxygen height $z_O$ at $c=9.9$~\AA{} (vacuum cell, three
collinear points: the production $z_O=0.96$~\AA{} and two dedicated checks at $z_O=0.90$ and
$1.02$~\AA) has slope $d\alpha/dz_O\approx-0.52$~eV/\AA$^2$ per \AA, the same charge-transfer
channel that sets $\Delta v$, made explicit as a single derivative rather than left implicit in the
total-energy scan. Only the slope is defensible here (three points define a line, not an absolute
stiffness); it is offered as the microscopic rate behind $\Delta v$, not as a fourth independent
calibration of it.

% ---- FIG S10: alpha vs oxygen height ----
\begin{figure}[t]
  \centering
  \figorbox{figS_zscan}{figS_zscan}{\columnwidth}
  \caption{The microscopic basis of the valence prediction, all vacuum (water-free) potentials at
  $c=9.9$~\AA. (a) The raw $E(\delta)$ curves at the three sampled oxygen heights $z_O=0.90$, 0.96
  (production), and $1.02$~\AA, shown for transparency. (b) The extracted Landau coefficient
  $\alpha(z_O)$: three collinear points give slope $d\alpha/dz_O\approx-0.52$~eV/\AA$^2$ per \AA. The
  defensible quantity is this slope, not an absolute well stiffness; it is the same
  oxygen-to-gallery charge-transfer channel that sets the valence split $\Delta v\approx0.06$ per
  cobalt quoted in the text.}
  \label{fig:zscan}
\end{figure}

\subsection{Gating and a negative control}

A LiCoO$_2$ thin film in which the carrier density is set electrostatically rather than chemically
should show the gallery turn-on without the disadvantage of the light frozen ion. Adding carriers to
the anhydrous (5.5~\AA) lithium cell by a compensating background, 0.1--0.3~$e$ per formula unit,
fills the gallery: the Fermi level and the areal density rise monotonically, to
$n_{\rm 2D}\approx4\times10^{14}$~cm$^{-2}$ at $0.3\,e$, and the lithium well softens as it fills
while remaining under a millielectronvolt deep, far below the lithium zero-point energy of
$\sim9$~meV, so the ion stays dynamically centred. Gating thus tunes the gas without freezing the
ion, exactly the place to test the gas and the magnetism apart from the glue, an ionic-liquid-gating
prediction for a thin-film collaboration.

Na$_2$CoSe$_2$O~\cite{Cheng2024}, a mixed-anion cobaltate ($T_c=6.3$~K) that carries the same
two-dimensional cobalt network with a different gallery chemistry, asks whether its gallery ion
hosts an interlayer band. It does not: the states at the Fermi level are the Co\,$3d$--Se\,$4p$
network, and the ionic sodium of its Na$_2$O gallery sits well above the Fermi level, with no
sodium-derived band there. Na$_2$CoSe$_2$O superconducts through its cobalt network, not the
gallery, and is a clean negative control: same CoO$_2$-like sheet, no gallery gas, superconductivity
by another route.

\bibliographystyle{apsrev4-2}
\bibliography{refs}

\end{document}
